% ============================================
% Johns Hopkins Letter Template
% Assistant Professor Cover Letter – Pomona College
% ============================================

\documentclass[11pt,letterpaper]{letter}

\usepackage{graphicx}
\usepackage{eso-pic}
\usepackage{geometry}
\usepackage{microtype}
\usepackage[hidelinks]{hyperref}

% ----- Page Setup -----
\geometry{left=25mm,right=25mm,top=25mm,bottom=25mm}

% ----- Logo Placement (first page only) -----
\newcommand{\PutJHULogoFG}[1]{%
  \AddToShipoutPictureFG*{%
    \AtPageUpperLeft{%
      \hspace{10mm}%
      \raisebox{-38mm}{%
        \includegraphics[width=6cm]{#1}%
      }%
    }%
  }%
}

% ----- Sender Information -----
\signature{Decory Edwards}
\address{Department of Economics \\ Johns Hopkins University \\ Baltimore, MD 21218 \\ \texttt{dedwar65@jh.edu}}

\begin{document}
\PutJHULogoFG{Images/jhulogo.png}

\begin{letter}{Search Committee \\
Department of Economics \\
Pomona College \\
425 N. College Ave \\
Claremont, CA 91711 \\ USA}

\opening{Dear Members of the Search Committee,}

I am writing to express my strong interest in the tenure-track Assistant Professor position in the Department of Economics at Pomona College. I am a Ph.D. candidate in Economics at Johns Hopkins University, advised by Professors Christopher D.~Carroll, Nicholas W.~Papageorge, and Stelios Fourakis, and I expect to receive my degree in May 2026. My research fields are macroeconomics, wealth and income dynamics, and computational economics.

My research agenda focuses on understanding the determinants of wealth inequality and the mechanisms through which heterogeneity in returns to wealth shapes aggregate outcomes. In my job-market paper, I estimate the degree of heterogeneity in individual portfolio returns using micro-data from the Survey of Consumer Finances and simulate a structural model in which households face idiosyncratic return risk. I show that allowing for persistent heterogeneity in returns substantially improves the model’s ability to match the empirical distribution of wealth, offering a tractable way to connect micro-level return dispersion to macroeconomic inequality.

In a second project, I use the Health and Retirement Study to examine the relationship between interpersonal trust and portfolio performance. Building on the literature that finds a hump-shaped relationship between trust and income, I show that a similar relationship holds for individual returns to wealth measured in the HRS data, highlighting a behavioral channel through which social attitudes shape saving and investment outcomes. This work also provides new estimates of the persistent component of returns to net worth, which motivate the calibration of heterogeneity in my structural model.

Pomona College offers an exceptional environment for combining rigorous research with deep student engagement. The 2–2 teaching load, strong institutional support for scholarship—including the fourth-year pre-tenure research leave—and the department’s effort to align teaching with faculty research interests make Pomona an especially compelling fit for my work. I am also excited by the opportunity to supervise senior theses, mentor students across the Claremont Colleges, and contribute to a residential liberal arts community committed to access, need-blind admissions, and supporting first-generation students. The collaborative ecosystem of the Claremont Colleges, with more than fifty economists across the consortium, is an ideal setting for my research in macroeconomics, inequality, and computational modeling.

\textbf{Teaching.} My teaching experience centers on undergraduate macroeconomics and an upper-level macro policy seminar, where I have led weekly discussion sections and held regular office hours for classes ranging from 16 to 40 students. My approach is student-centered: I aim to help students “convince themselves” that they have moved from confusion to understanding by holding structured office-hour coaching that builds trust and confidence. At Pomona, I am prepared to teach \textit{Principles of Macroeconomics}, \textit{Intermediate Macroeconomic Theory}, \textit{Money and Banking}, \textit{Econometrics}, and \textit{Financial Economics}, and I welcome the opportunity to supervise senior theses and mentor students across the Claremont Colleges.

I believe I would be a strong fit because my computational and empirical toolkit allows me to (i) produce rigorous, data-backed estimates of returns and distributional outcomes, (ii) build and validate quantitative models that speak to macroeconomic and policy-relevant questions, and (iii) mentor students effectively in a small liberal arts environment that values inclusive learning, high-impact research, and intellectual curiosity. Pomona’s balance of research excellence, close student interaction, and strong institutional support aligns closely with the academic community I hope to join.

I would be honored to contribute to Pomona College’s Department of Economics and to the broader Claremont Colleges community. Thank you for your consideration; I welcome the opportunity to discuss my fit with the department at your convenience.

\closing{Sincerely,}

\end{letter}
\end{document}
